% ============================================================
% HALAMAN RIWAYAT REVISI
% Letakkan sebelum \clearpage
\chapter*{ABSTRAK}
\addcontentsline{toc}{chapter}{ABSTRAK}

\begin{singlespace}
	\small
	Rekam Medis Elektronik (RME) merupakan komponen fundamental dalam penyelenggaraan pelayanan kesehatan di Indonesia. Peraturan Menteri Kesehatan Nomor 24 Tahun 2022 mewajibkan seluruh fasilitas pelayanan kesehatan menyelenggarakan RME dengan menjamin kerahasiaan, keutuhan, dan ketersediaan data. Untuk memenuhi kebutuhan tersebut, penelitian sebelumnya mengembangkan DecMed, sebuah kerangka kerja manajemen RME terdesentralisasi berbasis DLT IOTA, \textit{Proxy Re-Encryption} (PRE), dan IPFS. Meskipun DecMed berhasil menyelesaikan persoalan kontrol akses, mekanisme pencatatan aktivitasnya belum membentuk audit log yang memadai: pencatatan masih tersebar tanpa pengelola terpusat, cakupan \textit{event} sempit, catatan lama dapat tertimpa, format tidak terstandarisasi, dan bukti digital rentan terhapus akibat \textit{data pruning} pada jaringan DLT.

	Tugas Akhir ini mengembangkan \textit{Audit Log Service} (ALS), sebuah layanan \textit{backend} independen yang terintegrasi dengan arsitektur DecMed. Perancangan dimulai dengan analisis ancaman menggunakan metode STRIDE terhadap arsitektur DecMed, menghasilkan 25 ancaman yang dipetakan menjadi 13 jenis \textit{audit event} (EV1--EV13). Setiap \textit{event source} (Patient Client, Hospital Client, Ministry Client, dan PRE Server) mengirimkan \textit{audit event} ke ALS dalam format \textit{signed payload} Ed25519 secara asinkron. ALS memverifikasi tanda tangan digital, membentuk \textit{audit record} terstruktur, dan menerapkan \textit{hash-chaining} SHA-256 antarrecord untuk menjamin sifat \textit{tamper-evident} pada berkas log lokal. Secara periodik, ALS merotasi berkas log aktif, mengunggahnya ke IPFS Cluster dengan \textit{pinning} agar tidak terhapus oleh \textit{Garbage Collection}, lalu menerbitkan metadata rotasi sebagai \textit{frozen object} pada \textit{smart contract} Move di jaringan IOTA melalui Gas Station Server yang sudah tersedia pada DecMed.

	Hasil implementasi menunjukkan bahwa ALS berhasil mengagregasi \textit{audit event} dari seluruh komponen DecMed ke dalam satu repositori terpusat, merekam seluruh 13 jenis \textit{event} yang mencakup semua ancaman STRIDE yang teridentifikasi, mendeteksi perusakan rekam jejak melalui \textit{hash-chaining} kriptografis, serta menjamin retensi data jangka panjang melalui kombinasi \textit{IPFS pinning} dan metadata \textit{on-chain} yang \textit{immutable}. ALS tidak mengubah alur bisnis DecMed dan memanfaatkan kembali infrastruktur IPFS dan IOTA yang sudah tersedia.

	\textbf{\textit{Kata kunci: audit log, DecMed, IOTA, IPFS, hash-chaining, tanda tangan digital, STRIDE, rekam medis elektronik}}

\end{singlespace}
\clearpage
% ============================================================

\clearpage
\pagestyle{empty}

\begin{center}
    \Large\bfseries RIWAYAT REVISI LAPORAN
\end{center}

\vspace{0.4cm}

\begin{center}
    \small
    \begin{tabular}{ll}
        \textbf{Judul}      & : \textit{\thetitle} \\[4pt]
        \textbf{Penulis}    & : \theauthor \\[4pt]
        \textbf{NIM}        & : 13522162 \\[4pt]
        \textbf{Pembimbing} & : Dr.\ Phil.\ Eng.\ Hari Purnama, S.Si., M.Si \\
    \end{tabular}
\end{center}

\vspace{0.6cm}

\renewcommand{\arraystretch}{1.55}

\begin{xltabular}{\textwidth}{
    |>{\centering\arraybackslash}p{0.6cm}
    |X
    |>{\centering\arraybackslash}p{3.2cm}|}

\caption{Daftar Revisi Laporan Tugas Akhir}
\label{tab:riwayat-revisi}\\

\hline
\textbf{No} &
\textbf{Uraian Revisi} &
\textbf{Referensi} \\
\hline
\endfirsthead

\hline
\textbf{No} &
\textbf{Uraian Revisi} &
\textbf{Referensi} \\
\hline
\endhead

\hline
\endfoot

\hline
\endlastfoot

1 &
Menambah studi literatur terkait keilmuan security &
Subbab~\ref{sec:keamanan-informasi} \\
\hline

2 &
Menambah studi literatur terkait IOTA Address &
Subbab~\ref{subsec:karakteristik-audit-trail} \\
\hline

3 &
Menambah studi literatur terkait IPFS Pinning &
Tabel~\ref{tab:klasifikasi-ancaman} \\
\hline

4 &
Memindahkan Tabel Audit Event dari lampiran ke Badan Laporan  &
Tabel~\ref{tab:klasifikasi-ancaman} \\
\hline

5 &
Menambahkan Rancangan untuk mekanisme baru  &
Tabel~\ref{tab:klasifikasi-ancaman} \\
\hline

6 &
Memindahkan Tabel Audit Event dari lampiran ke Badan Laporan  &
Tabel~\ref{tab:klasifikasi-ancaman} \\
\hline

\end{xltabular}

\clearpage